\title{Comparing repeat measurements}
\subtitle{A practical guide to calibration}
\author{Methods seminar}
\institute{Teaching and Research Workshop}
\date{Reference edition}
\footlinepayoff{Shared reference \textperiodcentered{} Measurement methods}

\begin{document}
\maketitle
\section{Setup}

\begin{frame}{A repeat is only useful under known conditions}
Record the instrument, unit and measurement conditions.
\CoreParagraphToList
\begin{itemize}
\item Repeat the reading after resetting the instrument. Keep the conditions consistent so a change in the reading has a clear interpretation.
\item Keep the raw values alongside any average.
\item Flag a change of instrument or measurement unit.
\end{itemize}
\CoreBetweenGroups
\TermInline{Calibration}{Compare the instrument with a reference of known value.}
\end{frame}

\begin{frame}{A shift can survive repeated readings}
\TeachingQuestion{Do two nearly identical readings establish that the instrument is accurate?}
\CoreQuestionToNext
\CoreCaptioned{%
  \CoreTableLayout
  \begin{tabular}{@{}lll@{}}
    \toprule
    \textbf{Reference (mm)} & \textbf{First reading (mm)} & \textbf{Repeat (mm)}\\
    \midrule
    10.0 & 10.8 & 10.9\\
    15.0 & 15.8 & 15.9\\
    \bottomrule
  \end{tabular}}{Illustrative measurements with a common upward shift.}
\CoreTableToInterpretation
\TeachingAnswer{No. The readings agree closely with each other, but both exceed the reference value.

Check the calibration before interpreting those repeats. Even a stable instrument can report biased values when its baseline is wrong.}
\end{frame}

\section{Calibration procedures}
\begin{frame}{The calibration line shows systematic drift}
Compare the fitted line with the identity line.
\CoreParagraphToFigure
\CoreCaptioned{%
  \begin{tikzpicture}[x=1.5cm,y=1.45cm]
    \draw[coreInk] (0,0) -- (6,0) (0,0) -- (0,2.7);
    \draw[coreMuted,densely dashed] (.4,.24) -- (5.6,2.34);
    \draw[coreAccent,thick] (.4,.6) -- (5.6,2.62);
    \foreach \x/\y in {1/.83,2.1/1.32,3.5/1.79,4.8/2.36}
      \fill[coreAccent] (\x,\y) circle (1.6pt);
  \end{tikzpicture}}{Schematic plot of observed readings against reference values. The dashed line shows equality. The solid line and points illustrate a systematic upward shift.}
\CoreCaptionToBody
The gap matters even when the points cluster tightly around the fitted line.
\end{frame}

\begin{frame}{Two tests answer different questions}
\CoreEqualColumns{%
  \CoreSubhead{Repeatability}
  Read the same reference more than once.
  \CoreWithinGroup
  \begin{itemize}
  \item Look for variation within runs.
  \item Keep the reference fixed.
  \end{itemize}
}{%
  \CoreSubhead{Agreement}
  Compare readings with the known reference.
  \CoreWithinGroup
  \begin{itemize}
  \item Look for a persistent shift.
  \item Use several reference levels.
  \end{itemize}
}
\CoreBetweenGroups
\TeachingTakeaway{A repeatable instrument can still be biased.}
\end{frame}

\begin{frame}{One estimate needs an explanation}
\CoreEvidenceLeft{%
  \CoreSubhead{Paired readings}
  The same reference is measured before and after a reset.
  \CoreParagraphToDisplayMath
  \[
      \begin{aligned}
        \Delta_i &= y_{i,\mathrm{after}}-y_{i,\mathrm{before}},\\
        \mathcal{R} &= \{1,\dots,n\}.
      \end{aligned}
  \]
}{%
  \CoreSubhead{What the difference isolates}
  Pairing holds the reference fixed. Variation in the differences reflects a change in instrument response.
  \CoreWithinGroup
  \TermInline{Control}{Same reference and unit.}
}
\CoreBetweenGroups
\TeachingNote{Keep the individual differences available for inspection.}
\end{frame}

\begin{frame}{Report what was compared}
Use unbroken labels for measurement methods and units.
\CoreParagraphToTable
\begingroup\CoreTableLayout
\begin{tabularx}{\linewidth}{@{}p{.30\linewidth}YY@{}}
\toprule
\textbf{Method} & \textbf{What to record} & \textbf{What to check}\\
\midrule
\CoreAtom{Reference-standard} & Known reference and unit & Whether the reference covers the range\\
\CoreAtom{Paired-reset} & Readings before and after reset & Whether differences vary by reference\\
\bottomrule
\end{tabularx}
\CoreTableNote{Each row describes one comparison used in the calibration check.}
\endgroup
\CoreTableToInterpretation
The comparisons identify different ways a measurement can fail.
\end{frame}

\section{Apply}
\begin{frame}[fragile]{Compute the signed difference}
Subtract the reference reading. Keep the sign when you inspect the result.
\CoreWithinGroup
\begin{lstlisting}[language=Python]
def signed_error(readings, references):
    return [y - r for y, r in zip(readings, references)]
\end{lstlisting}
\CoreBetweenGroups
A positive difference means the reading exceeds the reference.
\end{frame}

\begin{frame}{Make the check repeatable}
Use the same sequence for each reference level.
\CoreParagraphToFigure
\begin{tikzpicture}[core process distance]
  \node[core process node] (load) {Set reference};
  \node[core process node,right=of load] (read) {Record reading};
  \node[core process node,right=of read] (reset) {Reset instrument};
  \node[core process node,right=of reset] (repeat) {Repeat reading};
  \draw[core process edge] (load) -- (read);
  \draw[core process edge] (read) -- (reset);
  \draw[core process edge] (reset) -- (repeat);
\end{tikzpicture}
\CoreBetweenGroups
\TeachingConclusion{A consistent reset lets you compare repeats at the same reference level.}
\end{frame}

\begin{frame}{Choose the next calibration check}
\slidefoot{Illustrative lab protocol}
\CoreEvidenceRight{%
  \CoreSubhead{If repeats vary}
  Compare readings immediately before and after the reset.
}{%
  \CoreSubhead{If the mean shifts}
  Compare several reference levels to see whether the error is constant or changes with the value.
}
\CoreBetweenGroups
\TeachingNote{In a bilingual report, use the same term for calibration (校准) throughout.}
\end{frame}

\CoreClosing{What a calibration report should explain}{
\begin{itemize}
\item State the reference values, units and measurement conditions.
\item Separate variation across repeated readings from bias against the reference.
\item Say which comparison would resolve the largest remaining uncertainty.
\end{itemize}
}
\end{document}
